%=====================================================
\section{Arc-length: Box 4.4 (cylindrical / Riks) formulation}
\label{sec:arclength_box44}

This section documents the cylindrical arc-length implementation in the
\texttt{solid\_mechanics} module corresponding to the input file
\texttt{arclength\_box44\_bar.i} and the verification script
\texttt{verify\_arclength\_box44\_bar.py}.

The algorithm is the \emph{non-consistent two-solve scheme}, Box~4.4 of
de~Souza~Neto, Peri\'{c}, Owen, \emph{Computational Methods for Plasticity}
(Wiley 2008), \S4.4.2.

%=====================================================
\subsection{Augmented system}

Standard quasi-static equilibrium with proportional loading reads
\begin{equation}
\bm{r}(\bu, \lambda) \;=\; \bm{f}_{\rm int}(\bu) \;-\; \lambda\, \hat{\bm{q}} \;=\; \bm{0},
\end{equation}
where $\bu \in \mathbb{R}^{n}$ collects the displacement degrees of freedom,
$\hat{\bm{q}} \in \mathbb{R}^{n}$ is the (fixed) reference external load, and
$\lambda \in \mathbb{R}$ is the load factor. We add the cylindrical arc-length
constraint
\begin{equation}
g(\Delta\bu) \;=\; \Delta\bu^{\top}\Delta\bu \;-\; \Delta s^2 \;=\; 0,
\qquad \Delta\bu = \bu - \bu_{n},
\end{equation}
where $\bu_{n}$ is the converged solution at the previous outer step and
$\Delta s$ is the user-prescribed arc-length increment.

%=====================================================
\subsection{Box 4.4 inner Newton iteration}

Within one outer arc-length step, the iteration $k=1,2,\dots$ proceeds as

\begin{enumerate}
\item Assemble $\bm{K}_T = \partial \bm{f}_{\rm int}/\partial \bu$ at $\bu^{(k-1)}$.

\item Two linear solves with the \emph{same} stiffness:
\begin{align}
\bm{K}_T\,\delta\bu^{*} &= -\bm{r}^{(k-1)}, \\
\bm{K}_T\,\delta\bar{\bu} &= +\hat{\bm{q}}.
\end{align}

\item Compute the iterative load increment $\delta\lambda^{(k)}$:
\begin{itemize}
\item For $k=1$ (\emph{predictor}):
\begin{equation}
\delta\lambda^{(1)} \;=\; \pm\,\frac{\Delta s}{\sqrt{\delta\bar{\bu}^{\top}\delta\bar{\bu}}},
\qquad
\text{sign} \;=\; \mathrm{sgn}\!\left(\Delta\bu_{\rm old}^{\top}\,\delta\bar{\bu}\right),
\end{equation}
defaulting to $+1$ on the first ever step. Here $\Delta\bu_{\rm old}$
is the previous \emph{converged} step's increment. Because
$\delta\bar{\bu} = \bm{K}_T^{-1}\hat{\bm{q}}$ is the elastic response to
the reference load, $\delta\bar{\bu}^{\top}\delta\bar{\bu} > 0$ by
construction, so the predictor has no first-step degeneracy.

\item For $k>1$ (\emph{corrector}): solve the scalar quadratic
\begin{equation}
a\,(\delta\lambda)^2 \;+\; b\,\delta\lambda \;+\; c \;=\; 0,
\end{equation}
with
\begin{align}
a &= \delta\bar{\bu}^{\top}\delta\bar{\bu}, \\
b &= 2\,(\Delta\bu^{(k-1)} + \delta\bu^{*})^{\top}\delta\bar{\bu}, \\
c &= (\Delta\bu^{(k-1)} + \delta\bu^{*})^{\top}(\Delta\bu^{(k-1)} + \delta\bu^{*}) - \Delta s^2.
\end{align}
Pick the root that maximises $\Delta\bu^{(k)}\!\cdot\!\Delta\bu^{(k-1)}$
(Crisfield 1991), where
\begin{equation}
\Delta\bu^{(k)} \;=\; \Delta\bu^{(k-1)} + \delta\bu^{*} + \delta\lambda\,\delta\bar{\bu}.
\end{equation}
\end{itemize}

\item Update load factor and displacement:
\begin{equation}
\lambda^{(k)} = \lambda^{(k-1)} + \delta\lambda^{(k)},
\qquad
\delta\bu^{(k)} = \delta\bu^{*} + \delta\lambda^{(k)}\,\delta\bar{\bu},
\qquad
\bu^{(k)} = \bu^{(k-1)} + \delta\bu^{(k)}.
\end{equation}

\item Recompute residual at $(\bu^{(k)}, \lambda^{(k)})$ and check
\begin{equation}
\frac{\|\bm{r}^{(k)}\|}{\max(|\lambda^{(k)}|, 1)\,\|\hat{\bm{q}}\|} \;<\; \varepsilon_r.
\end{equation}
\end{enumerate}

%=====================================================
\subsection{MOOSE implementation}

The algorithm is implemented entirely in MOOSE/libMesh, without any PETSc
SNES extensions. The architectural choices that proved necessary are:

\paragraph{$\lambda$ is an \texttt{AuxScalarVariable}, not a NL DOF.}
Putting $\lambda$ in the nonlinear system without a kernel filling its row
makes the assembled $\bm{K}$ singular at the $\lambda$ row, producing a
garbage $\delta\bar{\bu}$ from $\bm{K}_T\,\delta\bar{\bu}=\hat{\bm{q}}$.
Treating $\lambda$ as an externally controlled scalar — written into the
auxiliary system by \texttt{ArcLengthSolveObject} and read by the loading
BC via \texttt{coupledScalarValue} — keeps $\bm{K}_T$ a pure displacement
stiffness.

\paragraph{Reference load cached once.}
$\hat{\bm{q}}$ is obtained at the first call to \texttt{solve()} via
\begin{equation}
\hat{\bm{q}} \;=\; \bm{r}(\bu, \lambda{=}0) \;-\; \bm{r}(\bu, \lambda{=}1),
\end{equation}
which equals $\hat{\bm{q}}$ exactly under the proportional-loading
assumption used by the BC \texttt{ALCoupledScaledNeumannBC}. It is never
recomputed.

\paragraph{System update after solution writes.}
Both
\texttt{\_aux\_sys->system().update()} (after writing $\lambda$) and
\texttt{nl.system().update()} (after writing $\bu^{(k)}$) are required so
that variable reinit, which reads from \texttt{current\_local\_solution},
sees the new global state.

The objects that ship with the implementation are
\begin{center}
\begin{tabular}{ll}
\toprule
\texttt{ArcLengthSolveObject} & inner solve performing Box~4.4 \\
\texttt{ArcLengthTransient} & \texttt{Transient} executioner injecting the AL solve \\
\texttt{ALCoupledScaledNeumannBC} & contributes $-\lambda\,\hat{q}\,\phi_i$ \\
\bottomrule
\end{tabular}
\end{center}

%=====================================================
\subsection{Linear-elastic bar verification}
\label{subsec:arclength_box44_bar}

\paragraph{Domain and BCs.}
Plane-strain $10\times 1$ rectangle. The left edge is fully clamped:
$u_x = u_y = 0$ on $x=0$. The right edge ($x = L$) is loaded with the
arc-length-scaled Neumann traction
\begin{equation}
\bm{t}\big|_{x=L} \;=\; \lambda\,\sigma_{\rm ref}\,\bm{e}_x,
\qquad \sigma_{\rm ref} = 10^{3}~\text{MPa}.
\end{equation}

\paragraph{Material and parameters.}
\begin{table}[h]
\centering
\small
\begin{tabular}{lll}
\toprule
Symbol & Meaning & Value \\
\midrule
$E$ & Young's modulus & $10^{5}~\text{MPa}$ \\
$\nu$ & Poisson's ratio & $0.3$ \\
$L$ & bar length & $10~\text{mm}$ \\
$\sigma_{\rm ref}$ & reference end traction & $10^{3}~\text{MPa}$ \\
$\Delta s$ & arc-length step & $10^{-2}$ \\
$\varepsilon_r$ & relative residual tolerance & $10^{-6}$ \\
$N_{\rm steps}$ & number of outer steps & $10$ \\
\bottomrule
\end{tabular}
\end{table}

\paragraph{Closed form.}
For uniaxial stress in plane strain ($\sigma_{yy}=\sigma_{zz}=0$, $\varepsilon_{zz}=0$),
the effective modulus is
\begin{equation}
E_{\rm eff} \;=\; \frac{E}{1-\nu^{2}} \;=\; \frac{10^{5}}{0.91} \;\approx\; 1.0989\times 10^{5}~\text{MPa}.
\end{equation}
The tip displacement at load factor $\lambda$ is then
\begin{equation}
u_{\rm tip}^{\rm LEFM}(\lambda) \;=\; \lambda\,\frac{\sigma_{\rm ref}\,L}{E_{\rm eff}} \;=\; \lambda\cdot 0.0910~\text{mm}.
\label{eq:arclength_lefm}
\end{equation}

\paragraph{Result.}
The simulation traces a strictly linear $(\lambda, u_{\rm tip})$ trajectory.
A least-squares fit on the ten converged outer steps gives
\begin{equation}
\frac{u_{\rm tip}^{\rm sim}}{\lambda} \;=\; 0.09063~\text{mm},
\qquad
\frac{|0.09063 - 0.09100|}{0.09100} \;=\; 0.40\,\%,
\end{equation}
well inside the $0.5\,\%$ acceptance tolerance. Each outer step converges
in a single Newton iteration with $\|\bm{r}\|/\|\lambda\,\hat{\bm{q}}\|
\sim 10^{-15}$, consistent with linear elasticity reaching equilibrium
exactly at the predictor step.

\paragraph{Verification criterion.}
The pass criterion for this benchmark is
\begin{equation}
\frac{|m_{\rm sim} - m_{\rm LEFM}|}{m_{\rm LEFM}} \;<\; 5\times 10^{-3},
\qquad
m \;=\; \frac{u_{\rm tip}}{\lambda},
\end{equation}
which is satisfied with $0.40\,\%$ measured.
